\section{Vývojová omezení a proces}
\label{sec:kelvin-development}

Vývoj systému Kelvin je ovlivněn několika zásadními omezeními, která vyplývají jak z jeho historického vývoje, tak z organizačních podmínek jeho správy a údržby.
Tato omezení mají přímý dopad na návrh i samotnou implementaci řešení popsaného v této diplomové práci.

\subsection*{Atomicita změn a code review}

Jedním z hlavních pravidel vývoje je požadavek na malé a atomické pull requesty.
Každý pull request by měl obsahovat maximálně $\sim$~500 změněných řádků kódu a měl by se zaměřovat pouze na jednu konkrétní změnu nebo problém.
Cílem tohoto přístupu je zjednodušení procesu \uv{code review} a snížení rizika zavlečení chyb do produkčního systému.
V praxi to však znamená nutnost rozdělovat i logicky související úpravy do více menších kroků, což může celý vývojový proces prodlužovat.

\subsection*{Migrace frontendové části systému}

Dalším významným omezením je probíhající migrace frontendové části systému ze frameworku Svelte na framework Vue.
Při úpravách nebo přidávání nových funkcionalit by již neměly vznikat nové Svelte komponenty.
V případě rozsáhlejších změn je tak často nutné nejprve přepsat existující část aplikace do Vue a teprve následně provést požadovanou úpravu.
Tento postup výrazně zvyšuje časovou i technickou náročnost vývoje, zejména u funkcionalit zasahujících do uživatelského rozhraní.

Důvodem této migrace je zejména snaha o sjednocení technologického stacku, zlepšení dlouhodobé udržovatelnosti a snížení vstupní bariéry pro nové vývojáře systému.

\subsection*{Organizační omezení a plánování}

Proces schvalování změn představuje další omezení vývoje.
Code review je v současné době zajišťováno omezeným počtem vyučujících, kteří se systému věnují vedle svých dalších povinností.
Z tohoto důvodu se může stát, že zpracování pull requestu (žádost o sloučení změny) trvá i několik týdnů.
Tento fakt má přímý dopad na tempo vývoje a vyžaduje pečlivé plánování jednotlivých kroků implementace s dostatečným časovým předstihem.

\endinput
